% !TeX root = ../main.tex

\section{Semiconductor}

\begin{problem}[Elementary study of an intrinsic semiconductor]
  \begin{paracol}{2}
    Consider a semiconductor characterized by a constant density of states
    ($g(E) = C$) in both the valence and conduction band. These bands,
    separated by $E_g$, have a respective width of $E_V$ and $E_C$
    (see the Figure).
    \begin{enumext}
      \item Knowing that the Fermi level $E_F$ is in the bandgap
      ($5k_BT < E_F < E_g - 5k_BT$), leading to the conventional simplification
      for $f(E)$ (The Boltzmann approximation can be applied), find the
      expression for the number of electrons $n$ in the conduction band and the
      number of holes $h$ in the valence band respectively as a function of the
      data and absolute temperature.
    \end{enumext}
    \switchcolumn \centering
    \begin{tikzpicture}
      \draw [->] (0,0) node [left] {$0$} -- (3,0) node [right] {$g(E)$};
      \draw [->] (0,-2.8) -- (0,3.5) node [right] {$E$};
      \draw (0,-2.5) node [left] {$-E_V$} -- (1,-2.5) -- (1,0)
       node [above] {$C$};
      \draw (0,1) node [left] {$E_g$} -- (1,1) -- (1,3) -- (0,3)
       node [left] {$E_C + E_g$};
      \draw (0,.5) node [left] {$E_F$} --++ (.2,0);
    \end{tikzpicture}
  \end{paracol}
  \begin{enumext}[topsep = 0pt]
    \setcounter{enumXi}{1}
    \item Deduce the expression for the volumetric density of intrinsic
    carriers $n_i$ at ambient temperature ($k_BT = \qty{25}\meV$) as well as
    the position of the Fermi level $E_{F_i}$.

    Numerical applications: $E_g = \qty{0.7}\eV$, $E_C = E_V = \qty5\eV$,
    and $C = \qty{2e21}{electrons/\cm^3/\eV}$.
    \item At the same temperature, what is the intrinsic electronic
    conductivity of this material? Take
    $\mu_h = \mu_e = \qty{1000}{\cm^2/\V\s}$.
    \item The material is doped with $N_0$ boron impurities. Find the
    expressions and the new values of $n_d$, $h_d$, $\sigma_d$, and $E_{F_d}$
    with $N_0 = \num{e12}$ impurities per $\unit{\cm^3}$ and next with
    $N_0 = \num{e14}$ impurities per $\unit{\cm^3}$.
  \end{enumext}
\end{problem}
\begin{solution}\leavevmode
  \begin{enumext}
    \item The number of electrons in the conduction band
    \[
      n = \int_{E_g}^{E_g+E_c} g(E) f(E) \d E
        = Ck_BT\upe^{-\frac{(E_g-E_F)}{k_BT}}(1 - \upe^{-\frac{E_C}{k_BT}})
        \xlongequal[\text{Boltzmann approximation}]{E_C\gg k_BT}
          Ck_BT\upe^{-\frac{E_g-E_F}{k_BT}}.
    \]
    The number of electrons in the valence band
    \[
      h = \int_{-E_V}^0 g(E) [1 - f(E)] \d E
        \approxeq Ck_BT\upe^{-\frac{E_F}{k_BT}} (1 - \upe^{-\frac{E_V}{k_BT}})
        \xlongequal[\text{Boltzmann approximation}]{E_V\gg k_BT}
          Ck_BT\upe^{-\frac{E_F}{k_BT}}.
    \]
    where $f(E) \approxeq \upe^{-\frac{E-E_F}{k_BT}}$,
          $[1 - f(E)] \approxeq \upe^{\frac{E-E_F}{k_BT}}$.
    \item For the intrinsic carriers, the number of electrons and holes are the
    same
    \[
      n = Ck_BT\upe^{-\frac{E_g-E_F}{k_BT}} =
      h = Ck_BT\upe^{-\frac{E_F}{k_BT}}, \quad
      E_g - E_F = E_F.
    \]
    Then, the Fermi level is at mid-gap, i.e., $E_F = \frac{E_g}{2}$.
    Substitute it back into the expression of $n$ (or $h$),
    we get the intrinsic carrier concentration
    \[
      n_i = Ck_BT\upe^{-\frac{E_g}{2k_BT}}.
    \]
    Apply the parameters, $n_i = \qty{4.158e19}{\m^{-3}}$,
    $E_F = \qty{0.35}\eV$.
    \item The conductivity
    \[
      \sigma_i = e(n_i\mu_e + n_i\mu_h) = \qty{1.332}{\siemens/\m}.
    \]
    \item Boron is an acceptor, so doping concentration $N_A = N_0$.
    Now, the number of electrons and holes
    \[
      n = n_i,\ h = n_i \longrightarrow
      n_d = n,\ h_d = n + N_0,
    \]
    and they satisfy the identity: mass-action law
    \[
      n_i^2 = n_dh_d = n(n + N_0).
    \]
    We can solve that
    \[
      n = \frac{\sqrt{N_0^2 + 4n_i^2} - N_0}2,\quad
      n_d = n = \frac{\sqrt{N_0^2 + 4n_i^2} - N_0}2,\quad
      h = n + N_0 = \frac{\sqrt{N_0^2 + 4n_i^2} + N_0}2.
    \]
    The conductivity becomes
    \[
      \sigma_d = e(n_d\mu_e + h_d\mu_h).
    \]
    From $h = Ck_BT\upe^{-\frac{E_F}{k_BT}}$, we can obtain the Fermi level
    \[
      E_{F_d} = k_BT \ln\frac{C k_BT}{h_d}.
    \]
    Apply the parameters (with the previous $n_i = \qty{4.158e19}{\m^{-3}}$)
    \begin{enumext}[columns = 2]
      \item $N_0 = \qty{e18}{\m^{-3}}$.\\
      The numbers of the acceptors and donators
      \begin{gather*}
        n_d = \qty{4.108e19}{\m^{-3}},\\
        h_d = \qty{4.208e19}{\m^{-3}}.
      \end{gather*}
      The conductivity and the Fermi level
      \[
        \sigma_d = \qty{1.332}{\siemens/\m},~
        E_{F_d} = \qty{0.350}\eV.
      \]
      \item $N_0 = \qty{e20}{\m^{-3}}$.\\
      The numbers of the acceptors and donators
      \begin{gather*}
        n_d = \qty{1.503e19}{\m^{-3}},\\
        h_d = \qty{1.150e20}{\m^{-3}}.
      \end{gather*}
      The conductivity and the Fermi level
      \[
        \sigma_d = \qty{2.083}{\siemens/\m},~
        E_{F_d} = \qty{0.325}\eV.
      \]
    \end{enumext}
  \end{enumext}
\end{solution}